% AGH laboratory report. Engine: pdfLaTeX; bibliography: Biber.
% Content lives in content/, presentation in style/universityreport.sty.
\documentclass[11pt,a4paper,oneside]{article}
\usepackage{style/universityreport}
\addbibresource{references.bib}

% Required metadata. Keep plain text in title and author for PDF properties.
\reportinstitution{Akademia Górniczo-Hutnicza im. Stanisława Staszica w Krakowie}
\reportfaculty{Wydział Inżynierii Mechanicznej i Robotyki}
\reportprogramme{Inżynieria Akustyczna}
\reportcourse{Nazwa przedmiotu}
\reporttitle{Pomiar poziomu ciśnienia akustycznego}
\reportlabnumber{1}
\reportauthor{Imię Nazwisko, Imię Nazwisko}
\reportgroup{Grupa 1}
\reportinstructor{Tytuł, imię i nazwisko}
\reportdate{\today} % Date of exercise/submission: replace when needed.

% Optional metadata; empty values are hidden.
\reportlogo{figures/agh_small.jpg}
\reporttype{Sprawozdanie laboratoryjne} % Or: Raport techniczny
\reportsubtitle{}
\reportsemester{}
\reportmeta{}

\begin{document}
\makecoverpage
% For longer technical reports, uncomment the optional elements:
% % Optional: enable in main.tex for a longer technical report.
\section*{Streszczenie}
\addcontentsline{toc}{section}{Streszczenie}
Krótko opisać cel, zastosowaną metodę, najważniejszy wynik
oraz wniosek końcowy. Uzupełnić po napisaniu raportu.

% \tableofcontents

\section{Cel ćwiczenia}
\label{sec:cel}
Wyznaczyć średni poziom ciśnienia akustycznego w ustalonym punkcie
pomiarowym oraz oszacować niepewność wyniku.
\textbf{Wszystkie dane poniżej są demonstracyjne i należy je zastąpić
własnymi pomiarami.} Dostosować cel i zakres do instrukcji ćwiczenia.

\section{Podstawy teoretyczne}
\label{sec:teoria}
Poziom ciśnienia akustycznego opisuje zależność
\begin{equation}
  L_p = 20\log_{10}\!\left(\frac{p_{\mathrm{rms}}}{p_0}\right),
  \qquad p_0 = \qty{20}{\micro\pascal},
  \label{eq:lp}
\end{equation}
gdzie $p_{\mathrm{rms}}$ oznacza wartość skuteczną ciśnienia akustycznego.
Należy podać zastosowaną charakterystykę częstotliwościową i czasową
oraz warunki pomiaru. Przykład cytowania źródła: \autocite{lamport1994latex}.
W docelowym raporcie zastąpić je literaturą właściwą dla ćwiczenia.

\section{Aparatura i warunki pomiaru}
\label{sec:aparatura}
\begin{table}[htbp]
  \centering
  \caption{Aparatura pomiarowa -- pola do uzupełnienia.}
  \label{tab:aparatura}
  \begin{tabularx}{\linewidth}{@{}L L L@{}}
    \toprule
    Przyrząd & Model / numer & Zakres / dokładność \\
    \midrule
    Miernik poziomu dźwięku & Uzupełnić & Uzupełnić \\
    Kalibrator akustyczny & Uzupełnić & Uzupełnić \\
    \bottomrule
  \end{tabularx}
\end{table}
Opisać pomieszczenie, położenie źródła i mikrofonu, tło akustyczne,
czas pomiaru oraz ustawienia przyrządów. Podać warunki środowiskowe,
jeżeli wpływają na wynik.

\section{Metodyka}
\label{sec:metodyka}
\begin{enumerate}
  \item Sprawdzić tor pomiarowy zgodnie z instrukcją przyrządu.
  \item Ustalić geometrię stanowiska i zapisać ustawienia pomiarowe.
  \item Wykonać serię pomiarów w tych samych warunkach.
  \item Zapisać wyniki i sprawdzić wskazanie po zakończeniu serii.
\end{enumerate}
Opisać rzeczywisty przebieg ćwiczenia i odstępstwa od instrukcji.

\section{Wyniki pomiarów}
\label{sec:wyniki}
\Cref{tab:wyniki,fig:wyniki} przedstawiają tę samą serię danych
z pliku \texttt{data/measurements.csv}. Jest to przykład powtórzeń
w ustalonych warunkach, a nie poziomów przeznaczonych do sumowania.
\begin{table}[htbp]
  \centering
  \caption{Demonstracyjna seria pomiarów.}
  \label{tab:wyniki}
  \pgfplotstabletypeset[
    col sep=comma,
    columns={measurement,level},
    columns/measurement/.style={column name={Pomiar},column type=c},
    columns/level/.style={column name={\multicolumn{1}{c}{$L_p$ (\unit{\decibel})}},
      column type={S[table-format=2.1]},string type},
    every head row/.style={before row=\toprule,after row=\midrule},
    every last row/.style={after row=\bottomrule}
  ]{data/measurements.csv}
\end{table}
\begin{figure}[htbp]
  \centering
  \begin{tikzpicture}
    \begin{axis}[
      width=0.85\linewidth,height=45mm,
      xlabel={Numer pomiaru},ylabel={$L_p$ (\unit{\decibel})},
      xmin=0.5,xmax=5.5,ymin=69.5,ymax=70.5,xtick={1,2,3,4,5},
      grid=major,grid style={hairline},
      tick label style={font=\small},
      /pgf/number format/use comma,
    ]
      \addplot[accent,mark=*,thick] table[
        x=measurement,y=level,col sep=comma]{data/measurements.csv};
    \end{axis}
  \end{tikzpicture}
  \caption{Rozrzut demonstracyjnych wyników kolejnych powtórzeń.}
  \label{fig:wyniki}
\end{figure}

\section{Analiza niepewności}
\label{sec:niepewnosc}
Dla $n$ niezależnych powtórzeń tej samej wielkości, przy założeniu
stabilności warunków, średnią i jej niepewność typu A oblicza się jako
\begin{equation}
  \overline{L}_p = \frac{1}{n}\sum_{i=1}^{n} L_{p,i},
  \qquad
  u_A(\overline{L}_p) =
  \sqrt{\frac{\sum_{i=1}^{n}(L_{p,i}-\overline{L}_p)^2}{n(n-1)}}.
  \label{eq:niepewnosc}
\end{equation}
Dla danych demonstracyjnych $n=5$, $\overline{L}_p=\qty{70.0}{\decibel}$
i $u_A\approx\qty{0.07}{\decibel}$. Średnia arytmetyczna poziomów
opisuje tu estymację z powtórzeń; nie zastępuje uśredniania energetycznego
sygnałów zmiennych w czasie.

Uzupełnić budżet niepewności o składniki typu B wynikające z przyrządu,
kalibracji i metody. Przy niezależnych składnikach wyrażonych w jednostce
wyniku można zastosować
\begin{equation}
  u_c = \sqrt{u_A^2+\sum_j u_{B,j}^2},\qquad U=k u_c.
\end{equation}
Podać rozkłady, współczynniki wrażliwości i przyjęty współczynnik
rozszerzenia $k$. Nie przypisywać poziomu ufności bez uzasadnienia.

\section{Wnioski}
\label{sec:wnioski}
Podać wynik końcowy wraz z jednostką i niepewnością, odnieść go do celu
(zob. \cref{sec:cel}) i omówić najważniejsze ograniczenia metody.
Wskazać, co wynika z danych, oraz wyjaśnić ewentualne rozbieżności
z przewidywaniami. Zastąpić ten akapit własnymi wnioskami.

\printbibliography[heading=bibintoc,title={Literatura}]

% Optional supplementary material:
% \appendix
% % Optional: enable after \\appendix in main.tex.
\section{Materiały uzupełniające}
\label{sec:zalacznik}
Zamieścić dodatkowe dane, obliczenia lub kod niezbędny do odtworzenia analizy.
\begin{lstlisting}[language=Python,caption={Przykład obliczenia średniej i niepewności typu A.},label={lst:analiza}]
import csv
from statistics import mean, stdev
from math import sqrt

with open("data/measurements.csv", encoding="utf-8") as stream:
    levels = [float(row["level"]) for row in csv.DictReader(stream)]
print(mean(levels), stdev(levels) / sqrt(len(levels)))
\end{lstlisting}

\end{document}
